% generated by analysis.py; do not edit
\newcommand{\SineCAGR}{16.17\%}
\newcommand{\SineVol}{12.07\%}
\newcommand{\SineSharpe}{1.30}
\newcommand{\SineMaxDD}{-11.85\%}
\newcommand{\LinearCAGR}{16.16\%}
\newcommand{\LinearVol}{12.06\%}
\newcommand{\LinearSharpe}{1.30}
\newcommand{\LinearMaxDD}{-11.87\%}
\newcommand{\PolynomialCAGR}{16.04\%}
\newcommand{\PolynomialVol}{12.04\%}
\newcommand{\PolynomialSharpe}{1.30}
\newcommand{\PolynomialMaxDD}{-12.12\%}
\newcommand{\HedgedCAGR}{15.92\%}
\newcommand{\HedgedVol}{11.24\%}
\newcommand{\HedgedSharpe}{1.37}
\newcommand{\HedgedMaxDD}{-10.81\%}
\newcommand{\CoreCAGR}{12.32\%}
\newcommand{\CoreVol}{11.19\%}
\newcommand{\CoreSharpe}{1.09}
\newcommand{\CoreMaxDD}{-12.72\%}
\newcommand{\SineTurnover}{3.91\%}
\newcommand{\LinearTurnover}{3.93\%}
\newcommand{\PolynomialTurnover}{4.04\%}
\newcommand{\HOneResult}{Not supported}
\newcommand{\HTwoResult}{Supported}
\newcommand{\HThreeResult}{Supported}
\newcommand{\LatestModelDate}{2026-09-18}
\newcommand{\LatestModelCoordinate}{0.00}
\newcommand{\LatestSPY}{57.00\%}
\newcommand{\LatestTLT}{38.00\%}
\newcommand{\LatestGLD}{3.50\%}
\newcommand{\LatestSLV}{1.50\%}


\section{Thesis-specific model}\label{sec:thesis-math}

GLD seeks to reflect gold-bullion prices less trust expenses, while SLV seeks to track silver
bullion. They are liquid securities and imperfect commodity proxies, not direct delivery contracts.
To control estimation error, the four assets are reduced to two internally fixed sleeves. Their
returns are
\begin{equation}\label{eq:tm-sleeves}
 c_t=0.60r_{\mathrm{SPY},t}+0.40r_{\mathrm{TLT},t},
 \qquad
 m_t=0.70r_{\mathrm{GLD},t}+0.30r_{\mathrm{SLV},t}.
\end{equation}
For a metals allocation \(h\in[0.05,0.35]\), the four-asset weight vector is
\begin{equation}\label{eq:tm-weights}
 w(h)=\left(0.60(1-h),\ 0.40(1-h),\ 0.70h,\ 0.30h\right)^{\top}.
\end{equation}
The static hedged reference has \(h_0=0.20\), hence weights
\((0.48,0.32,0.14,0.06)^{\top}\).

Let \(x_t=(c_t,m_t)^{\top}\). Using at most \(L=126\) completed sessions, the rolling estimates
are
\begin{equation}\label{eq:tm-moments}
 \widehat\mu_t=252\,\overline x_t,
 \qquad
 \widehat\Sigma_t=252\,\widehat{\operatorname{Cov}}_t(x)+\epsilon I.
\end{equation}
The global minimum-variance metals share and the mean--variance growth anchor are
\begin{equation}\label{eq:tm-anchor}
 h_t^{\mathrm{gmv}}=\arg\min_h v(h)^{\top}\widehat\Sigma_t v(h),
 \quad
 h_t^{\star}=\arg\max_{h\in\mathcal F_t}
 \left\{v(h)^{\top}\widehat\mu_t-\frac{\gamma}{2}
 v(h)^{\top}\widehat\Sigma_t v(h)\right\},
 \quad v(h)=(1-h,h)^{\top},
\end{equation}
where \(\gamma=3\) and \(\mathcal F_t\) is the estimated efficient segment from the GMV point to
the higher-estimated-return endpoint.

The five-session compound-growth gap between the anchor and static hedged reference is
\begin{equation}\label{eq:tm-growth-gap}
 d_t=\prod_{j=0}^{4}\left(1+w(h_t^{\star})^{\top}r_{t-j}\right)
 -\prod_{j=0}^{4}\left(1+w(h_0)^{\top}r_{t-j}\right).
\end{equation}
If \(\widehat\sigma_{5,t}\) and \(\widehat\sigma_{21,t}\) are annualised volatilities of the
static hedged reference, the raw coordinate is
\begin{equation}\label{eq:tm-coordinate}
 q_t=\operatorname{clip}\!\left(
 \frac{\widehat\sigma_{5,t}/\widehat\sigma_{21,t}-1}{0.50},-1,1\right)
 \min\!\left\{\frac{|d_t|}{\widehat\sigma_{21,t}\sqrt{5/252}},1\right\}.
\end{equation}
Positive \(q_t\) moves toward the lower-risk point; negative \(q_t\) moves toward the frontier's
higher-risk endpoint. A zero growth gap leaves the allocation at the rolling anchor.

For interpolation type \(k\), the causal five-session coordinate is
\begin{equation}\label{eq:tm-interpolation}
 \widetilde q_t^{(k)}=\sum_{j=1}^{5}\kappa_j^{(k)}q_{t-5+j},
 \qquad
 \kappa_j^{(k)}=\frac{g_k(j)}{\sum_{\ell=1}^{5}g_k(\ell)},
\end{equation}
with \(g_{\mathrm{sine}}(j)=\sin(\pi j/10)\), \(g_{\mathrm{linear}}(j)=j\), and
\(g_{\mathrm{poly}}(j)=j^2\). The target formed after close \(t\) is applied to the return on
session \(t+1\).

\section{Empirical comparison}\label{sec:interpolation-comparison}

Table~\ref{tab:interpolations} separates two questions: whether adding a static metals sleeve
helped relative to unhedged 60/40, and whether dynamically timing that sleeve added value.

\begin{table}[H]
\centering
\caption{Net performance from January 2023.}
\label{tab:interpolations}
\begin{tabular}{@{}lrrrrr@{}}
\toprule
Rule & CAGR & Volatility & Sharpe & Max drawdown & Avg. turnover \\
\midrule
Sine & \SineCAGR{} & \SineVol{} & \SineSharpe{} & \SineMaxDD{} & \SineTurnover{} \\
Linear & \LinearCAGR{} & \LinearVol{} & \LinearSharpe{} & \LinearMaxDD{} & \LinearTurnover{} \\
Quadratic & \PolynomialCAGR{} & \PolynomialVol{} & \PolynomialSharpe{} & \PolynomialMaxDD{} & \PolynomialTurnover{} \\
Static hedged & \HedgedCAGR{} & \HedgedVol{} & \HedgedSharpe{} & \HedgedMaxDD{} & -- \\
Static 60/40 & \CoreCAGR{} & \CoreVol{} & \CoreSharpe{} & \CoreMaxDD{} & -- \\
\bottomrule
\end{tabular}
\end{table}

The pre-registered outcomes are H1 \textbf{\HOneResult{}}, H2 \textbf{\HTwoResult{}}, and H3
\textbf{\HThreeResult{}}. These labels describe one historical sample, not a population result.

\begin{figure}[H]
\centering
\includegraphics[width=0.96\textwidth]{figures/fig_metals_wealth.pdf}
\caption{Top: cumulative wealth for the primary sine rule, static references, and SPY. Bottom:
linear and quadratic cumulative-wealth differences from sine, in basis points.}
\label{fig:metals-wealth}
\end{figure}

\begin{figure}[H]
\centering
\includegraphics[width=0.96\textwidth]{figures/fig_metals_weights.pdf}
\caption{Top: sine-rule metals target and the static 20\% reference. Bottom: target-weight
differences between the alternative kernels and sine.}
\label{fig:metals-weights}
\end{figure}

\begin{figure}[H]
\centering
\includegraphics[width=0.96\textwidth]{figures/fig_metals_signal.pdf}
\caption{Reference volatility estimates and the sine-smoothed frontier coordinate; the noisy raw
coordinate is shown faintly for context.}
\label{fig:metals-signal}
\end{figure}

\begin{figure}[H]
\centering
\includegraphics[width=0.90\textwidth]{figures/fig_metals_hedge_drawdown.pdf}
\caption{Drawdown paths for the static metals-hedged reference and unhedged 60/40.}
\label{fig:metals-drawdown}
\end{figure}

\begin{figure}[H]
\centering
\includegraphics[width=0.86\textwidth]{figures/fig_metals_frontier.pdf}
\caption{Latest estimated frontier between the fixed core and metals sleeves. Coincident GMV,
anchor, and target points are collapsed into one labelled marker.}
\label{fig:metals-frontier}
\end{figure}

\begin{figure}[H]
\centering
\includegraphics[width=0.86\textwidth]{figures/fig_metals_costs.pdf}
\caption{Top: interpolation Sharpe ratios under alternative trading costs. Bottom: linear and
quadratic differences from sine, magnified by 100.}
\label{fig:metals-costs}
\end{figure}

\section{Trade and broker interpretation}\label{sec:broker-interpretation}

The latest completed-session target, dated \LatestModelDate{}, is \LatestSPY{} SPY,
\LatestTLT{} TLT, \LatestGLD{} GLD, and \LatestSLV{} SLV. Its smoothed coordinate is
\LatestModelCoordinate{}. These are model weights, not an executable recommendation. The supplied
IBKR planner checks for paper mode, reads current positions and base-currency prices, sizes
whole-share differences, and sets \texttt{transmit=False}; it never places an order.

\section{References}\label{sec:references}

\begin{itemize}
  \item H. Markowitz (1952), ``Portfolio Selection,'' \emph{Journal of Finance} 7(1), 77--91,
  \href{https://doi.org/10.1111/j.1540-6261.1952.tb01525.x}{doi:10.1111/j.1540-6261.1952.tb01525.x}.
  \item World Gold Trust Services, ``SPDR Gold Shares: GLD,''
  \href{https://www.spdrgoldshares.com/usa/gld/}{official fund page}.
  \item iShares, ``iShares Silver Trust: SLV,''
  \href{https://www.ishares.com/us/products/239855/ishares-silver-trust-fund}{official fund page}.
  \item D. N. Politis and J. P. Romano (1994), ``The Stationary Bootstrap,'' \emph{JASA} 89(428),
  1303--1313, \href{https://doi.org/10.1080/01621459.1994.10476870}{doi:10.1080/01621459.1994.10476870}.
\end{itemize}
